% !TEX root = ../main.tex
% =====================================================================
% ABSTRACT  (~1/2 page, one paragraph, no citations, no figures)
% Write this LAST, once every number is frozen.
% =====================================================================

% ---------------------------------------------------------------------
% WRITING GUIDE (as comments, because the abstract is a quote environment
% and a grey guide box cannot be placed inside one).
%
% One paragraph, about 150-200 words. Recipe, one or two sentences each:
%   1. THE PROBLEM      - deep models classify time series well but do not say
%                         WHICH PART of the signal caused the decision.
%   2. EXISTING ANSWER  - MILLET makes the model explain itself with MIL pooling.
%   3. WHAT I BUILT     - SEA-Net: a small multi-scale encoder plus my own family
%                         of MIL pooling heads.
%   4. HOW I TESTED IT  - WebTraffic + the 85 UCR datasets, 3 seeds, accuracy and
%                         explanation quality (AOPCR, NDCG@n).
%   5. HEADLINE NUMBER  - one clear result: accuracy ~ baseline, explanations
%                         better, parameters down.
%
% Questions to answer before writing:
%   - If a reader only reads this paragraph, what single sentence do I want them
%     to remember?
%   - What is the one number that proves it?
% ---------------------------------------------------------------------


In tiny deep learning, previous models were black boxes. They only predicted classes or assigned labels to time series data without providing any reason for it. The \millet{} framework showed that Multiple Instance Learning (MIL) pooling turns a standard backbone into a model that explains itself. The network scores every time step, and those scores \emph{are} the explanation. However, the issue is that this model was not made for tiny applications.

During this internship, I built \seanet{}, which is our baseline architecture. In this model, we use a Temporal Convolutional Network (TCN) as the main encoder along with different pooling techniques. My main contributions are the multi-scale depthwise-separable TCNs that create the \seanet{} family of encoders (such as \texttt{sea\_mstcn\_sep}, \texttt{sea\_mstcn\_sep\_gated}, \texttt{sea\_mstcn\_sep\_bottleneck}, \texttt{sea\_mstcn\_sep\_inputgate}, and \texttt{sea\_mstcn\_sep\_recon}). I also contributed to the pooling family, including \texttt{classwise\_conjunctive}, \texttt{softmax\_conjunctive}, \texttt{adaptive\_classwise\_topk\_conjunctive}, \texttt{attention\_max}, \texttt{gated\_attention}, and \texttt{dualstream\_conjunctive}.

I evaluated \num{67} model combinations on the \webtraffic{} dataset and the 85 \ucr{} datasets used by \millet{}. I measured both classification quality (accuracy, loss) and explanation quality (AOPCR, NDCG@$n$) using one random seed. However, to compare with the \millet{} baseline, we evaluated our top model, \texttt{seanet\_bottleneck\_topk}, over \num{3} random seeds. My best variant reaches 95\% accuracy (\num{0.95}) on \webtraffic{}, compared to \num{0.887} for the re-trained \millet{} baseline. At the same time, it uses \num{10}$\times$ fewer parameters and improves the NDCG score (used to evaluate explanations) from \num{0.677} to \num{0.772}.
